%Data institusi

%\input{Dozentdaten}

%\renewcommand{\fachbereich}{Fachbereich}

%\renewcommand{\dozent}{Prof. Dr. . }

%Data ujian

\renewcommand{\klausurgebiet}{ }

\renewcommand{\klausurtyp}{ }

\renewcommand{\klausurdatum}{ . 20}

\klausurvorspann {\fachbereich} {\klausurdatum} {\dozent} {\klausurgebiet} {\klausurtyp}

%Data untuk tabel nilai berikut

\renewcommand{\aeins}{  3  }

\renewcommand{\azwei}{  3   }

\renewcommand{\adrei}{  0  }

\renewcommand{\avier}{  0  }

\renewcommand{\afuenf}{  6  }

\renewcommand{\asechs}{  6  }

\renewcommand{\asieben}{  0  }

\renewcommand{\aacht}{  6  }

\renewcommand{\aneun}{  11  }

\renewcommand{\azehn}{  5  }

\renewcommand{\aelf}{  5  }

\renewcommand{\azwoelf}{  3  }

\renewcommand{\adreizehn}{  7   }

\renewcommand{\avierzehn}{  2  }

\renewcommand{\afuenfzehn}{  0  }

\renewcommand{\asechzehn}{  6  }

\renewcommand{\asiebzehn}{  63  }

\renewcommand{\aachtzehn}{    }

\renewcommand{\aneunzehn}{    }

\renewcommand{\azwanzig}{    }

\renewcommand{\aeinundzwanzig}{    }

\renewcommand{\azweiundzwanzig}{    }

\renewcommand{\adreiundzwanzig}{    }

\renewcommand{\avierundzwanzig}{    }

\renewcommand{\afuenfundzwanzig}{    }

\renewcommand{\asechsundzwanzig}{    }

\punktetabellesechzehn

\klausurnote

\newpage

\setcounter{section}{K}



\inputaufgabeklausurloesung
{3}

{

Definisikan istilah-istilah berikut
\zusatzklammer {yang dicetak miring} {} {.}
\aufzaehlungsechs{Suatu
\stichwort {kurva geodesik} {}
pada hiperpermukaan diferensiabel

\mavergleichskette
{\vergleichskette
{ Y
}
{ \subseteq }{ \R^n
}
{  }{
}
{    }{
}
{   }{
}
}
{}{}{.}

}{
\stichwort {Himpunan rotasi} {}
\zusatzklammer {mengelilingi sumbu $x$} {} {}
dari suatu subhimpunan

\mavergleichskette
{\vergleichskette
{ T
}
{ \subseteq }{  \R \times \R_{\geq 0}
}
{  }{
}
{    }{
}
{   }{
}
}
{}{}{.}

}{
\stichwort {Pemetaan tangen} {}
\maabbdisp {T (\varphi)} {TM} {TN
} {}
dari
\definitionsverweis {pemetaan diferensiabel}{}{}
\maabbdisp {\varphi} { M } { N
} {}
antara
\definitionsverweis {manifold-manifold diferensiabel}{}{}
\mathkor {} {M} {dan} {N} {.}

}{Suatu perubahan peta yang
\stichwort {mempertahankan orientasi} {}
pada
\definitionsverweis {manifold diferensiabel}{}{}
$M$.

}{
\stichwort {Bentuk diferensial tarikan balik} {} $\varphi^* \omega$ dari suatu bentuk diferensial
\mathl{\omega \in
{ \mathcal E }^{ k } (  M   )}{} terhadap suatu pemetaan terdiferensialkan secara kontinu
\maabb {\varphi} { L } { M
} {}
antara dua manifold diferensiabel
\mathkor {} {L} {dan} {M} {.}

}{Suatu koneksi linear yang
\stichwort {metrik} {}
pada bundel tangen suatu manifold Riemann $M$.
}

}
{

\aufzaehlungsechs{Suatu
\zusatzklammer {sebagai pemetaan ke $\R^n$} {} {}
\definitionsverweis {kurva yang terdiferensialkan secara kontinu dua kali}{}{}
\maabbdisp {\gamma} {I} {Y
} {}
disebut
kurva geodesik
pada $Y$ jika
\definitionsverweis {percepatan tangensial}{}{}
kurva tersebut sama dengan $0$ di setiap titik.
}{Himpunan rotasi dari $T$ adalah

\mathdisp {\left\{ (x, y \cos \alpha, y \sin \alpha  ) \in \R^3  \mid   (x,y) \in T , \, \alpha \in [0, 2 \pi]       \right\}} { . }

}{Yang dimaksud dengan pemetaan tangen
\maabbdisp {T (\varphi)} {TM} {TN
} {}
adalah gabungan saling lepas dari
\definitionsverweis {pemetaan-pemetaan tangen}{}{}
di setiap titik, yakni

\mavergleichskettedisp
{\vergleichskette
{ T (\varphi)
}
{ =} { \biguplus_{P \in M} T_P(\varphi)
}
{ } {
}
{ } {
}
{ } {
}
}
{}{}{.}
}{Misalkan
\maabbdisp {\alpha_1} {U_1 } { V_1
} {}
dan
\maabbdisp {\alpha_2} {U_2 } { V_2
} {}
adalah peta-peta berorientasi dari $M$. Perubahan peta terkait
\definitionsverweis {perubahan peta}{}{}
\maabbdisp {\psi=\alpha_2 \circ \alpha_1^{-1}} {V_1 \cap \alpha_1 (U_1 \cap U_2)} {V_2 \cap \alpha_2 (U_1 \cap U_2 )
} {}
disebut mempertahankan orientasi jika untuk setiap titik
\mathl{Q \in V_1 \cap \alpha_1 (U_1 \cap U_2)}{},
\definitionsverweis {diferensial total}{}{}
\maabbdisp {\left(D\psi\right)_{Q}} {\R^n } {\R^n
} {}
\definitionsverweis {mempertahankan orientasi}{}{}.
}{Bentuk diferensial tarikan balik
\mathl{\varphi^* \omega}{} didefinisikan, untuk
\mathl{P\in L}{} dan
\mathl{v_1 , \ldots , v_k \in T_PL}{}, oleh

\mavergleichskettedisp
{\vergleichskette
{  ( \varphi^* \omega )(P,v_1 , \ldots , v_k )
}
{ \defeq} { \omega(\varphi(P), T_P\varphi(v_1) , \ldots ,  T_P\varphi(v_k) )
}
{ } {
}
{ } {
}
{ } {
}
}
{}{}{}
.
}{Koneksi tersebut disebut
metrik
jika

\mavergleichskettedisp
{\vergleichskette
{ D_V \left\langle W , Z  \right\rangle
}
{ =} { \left\langle \nabla_V W , Z  \right\rangle  + \left\langle W ,  \nabla_VZ  \right\rangle
}
{ } {
}
{ } {
}
{ } {
}
}
{}{}{.}
berlaku untuk sebarang medan vektor $V,W ,Z$.
}

 }



\inputaufgabeklausurloesung
{3}

{

Nyatakan teorema-teorema berikut.
\aufzaehlungdrei{
\stichwort {Teorema tentang solusi persamaan diferensial biasa pada hiperpermukaan diferensiabel} {}

\mavergleichskette
{\vergleichskette
{  Y
}
{ \subseteq }{  \R^n
}
{  }{
}
{    }{
}
{   }{
}
}
{}{}{.}}{Rumus untuk menghitung volume kanonik suatu manifold Riemann dalam sebuah peta.}{Teorema tentang partisi kesatuan.}

}
{

\aufzaehlungdrei{\faktsituation {Misalkan

\mavergleichskette
{\vergleichskette
{ W
}
{ \subseteq }{ \R^n
}
{  }{
}
{    }{
}
{   }{
}
}
{}{}{}
\definitionsverweis {terbuka}{}{,}
\maabb {h} { W } { \R
} {}
suatu
\definitionsverweis {fungsi yang terdiferensialkan secara kontinu}{}{}
dan

\mavergleichskette
{\vergleichskette
{ Y
}
{ = }{ h^{-1}(c)
}
{  }{
}
{    }{
}
{   }{
}
}
{}{}{}
\definitionsverweis {serat}{}{}
atas

\mavergleichskette
{\vergleichskette
{ c
}
{ \in }{ \R
}
{  }{
}
{    }{
}
{   }{
}
}
{}{}{,}
dengan $h$ bersifat
\definitionsverweis {reguler}{}{}
di setiap titik $Y$.
Misalkan

\mavergleichskette
{\vergleichskette
{ Y
}
{ \subseteq }{ U
}
{ \subseteq }{ \R^n
}
{    }{
}
{   }{
}
}
{}{}{}
\definitionsverweis {terbuka}{}{}
dan misalkan $F$ suatu
\definitionsverweis {medan vektor}{}{}
terdiferensialkan pada $U$}
\faktvoraussetzung {dengan sifat bahwa untuk setiap titik

\mavergleichskette
{\vergleichskette
{ P
}
{ \in }{ Y
}
{  }{
}
{    }{
}
{   }{
}
}
{}{}{}
vektor $F(P)$ termasuk dalam
\definitionsverweis {ruang tangen serat}{}{}
di $P$ pada $Y$.}
\faktfolgerung {Maka solusi $v(t)$ dari
\definitionsverweis {masalah nilai awal}{}{}
untuk $F$ dengan syarat awal

\mavergleichskette
{\vergleichskette
{ v(t_0)
}
{ \in }{ Y
}
{  }{
}
{    }{
}
{   }{
}
}
{}{}{}
seluruhnya terletak di $Y$.}
\faktzusatz {}
\faktzusatz {}}{Misalkan $M$ suatu
\definitionsverweis {manifold Riemann}{}{}
\definitionsverweis {berorientasi}{}{}
dan $\omega$ adalah
\definitionsverweis {bentuk volume kanonik}{}{.}
Misalkan
\maabbdisp {\alpha} { U } { V
} {}
suatu
\definitionsverweis {peta berorientasi}{}{}
dengan

\mavergleichskettedisp
{\vergleichskette
{ V
}
{ \subseteq} { \R^n
}
{ } {
}
{ } {
}
{ } {
}
}
{}{}{}
terbuka, berkoordinat
\mathl{x_1 , \ldots , x_n}{}, dengan
\definitionsverweis {matriks fundamental metrik}{}{}

\mathl{G=(g_{ij})_{1 \leq i,j \leq n}}{} dan
\mathl{g= \det  G}{.} Maka

\mavergleichskettedisp
{\vergleichskette
{  \alpha_* \omega
}
{ =} { \sqrt{ g} dx_1 \wedge \ldots \wedge dx_n
}
{ } {
}
{ } {
}
{ } {
}
}
{}{}{.}
Untuk suatu
\definitionsverweis {subhimpunan terukur}{}{}

\mavergleichskette
{\vergleichskette
{ T
}
{ \subseteq }{ U
}
{  }{
}
{    }{
}
{   }{
}
}
{}{}{}
berlaku

\mavergleichskettedisp
{\vergleichskette
{
\int_{  T   }  \omega
}
{ =} {
\int_{ \alpha(T)  }   \sqrt{g} dx_1 \wedge \ldots \wedge dx_n
}
{ =} { \int_{ \alpha(T)  }   \sqrt{g}    \,  d \lambda^n
}
{ } {
}
{ } {
}
}
{}{}{.}}{Misalkan $M$ suatu
\definitionsverweis {manifold diferensiabel}{}{}
dengan
\definitionsverweis {basis terhitung}{}{}
bagi topologinya. Maka untuk setiap penutup terbuka terdapat
\definitionsverweis {partisi kesatuan}{}{}
yang
\definitionsverweis {terdiferensialkan secara kontinu}{}{}
dan subordinat terhadap penutup tersebut.}

 }



\inputaufgabeklausurloesung
{0}

{

}
{  }



\inputaufgabeklausurloesung
{0}

{

}
{  }



\inputaufgabeklausurloesung
{6}

{

Untuk permukaan $Y$ yang diberikan oleh

\mavergleichskettedisp
{\vergleichskette
{ x^2+y^2+2z^2
}
{ =} { 4
}
{ } {
}
{ } {
}
{ } {
}
}
{}{}{}
dan titik

\mavergleichskettedisp
{\vergleichskette
{ P
}
{ =} { \left(  1  , \,   1  , \,   1                 \right)
}
{ \in} { Y
}
{ } {
}
{ } {
}
}
{}{}{,}
tentukan suatu matriks diagonal bagi
\definitionsverweis {pemetaan Weingarten}{}{}
$L_P$.

}
{

Medan gradien ternormalisasi adalah

\mavergleichskettealign
{\vergleichskettealign
{ N(x,y,z)
}
{ =} { { \frac{   1  }{  \sqrt{ 4x^2+4y^2+16 z^2}  } } \begin{pmatrix}  2x \\ 2y\\  4z   \end{pmatrix}
}
{ =} { { \frac{   1  }{  \sqrt{ x^2+y^2+4 z^2}  } } \begin{pmatrix}  x  \\ y \\  2z   \end{pmatrix}
}
{ } {
}
{ } {
}
}
{}
{}{.}
Kita bekerja dengan medan normal satuan

\mavergleichskettedisp
{\vergleichskette
{ M(x,y,z)
}
{ =} { { \frac{   1  }{  \sqrt{ 4+2 z^2}  } } \begin{pmatrix}  x  \\ y \\  2z   \end{pmatrix}
}
{ } {
}
{ } {
}
{ } {
}
}
{}{}{,}
yang pada $Y$ berimpit dengan $N$. Diferensial total dari $M$ adalah

\mathdisp {\begin{pmatrix}  { \frac{   1  }{  \sqrt{4+2z^2}  } }   &   0 &  { \frac{   -2 xz }{  \left(  4+2z^2  \right)^{3/2}  } }  \\  0  &  { \frac{   1  }{  \sqrt{4+2z^2}  } }   &  { \frac{   -2 yz }{  \left(  4+2z^2  \right)^{3/2}  } }   \\ 0   &  0  &  { \frac{   8 }{  \left(  4+2z^2  \right)^{3/2}  } }  \end{pmatrix}} { . }

Di titik $P$ yang diberikan, gradiennya adalah $\begin{pmatrix}  2  \\ 2\\  4   \end{pmatrix}$ dan kedua vektor
\mathkor {} {\begin{pmatrix}  1  \\ -1\\  0   \end{pmatrix}} {dan} {\begin{pmatrix}  2  \\ 0 \\  -1   \end{pmatrix}} {}
membentuk suatu basis ruang tangen $T_PY$. Diferensial total dari $M$ di titik ini sama dengan

\mathdisp {\begin{pmatrix}  { \frac{   1  }{  \sqrt{6}  } }   &   0 &  { \frac{   -1 }{  3 \sqrt{6}  } }  \\  0  &  { \frac{   1  }{  \sqrt{6}  } }   &  { \frac{   -1  }{  3 \sqrt{6}  } }   \\ 0   &  0  &  { \frac{   4  }{  3 \sqrt{6}  } }   \end{pmatrix}} { . }

Jika diterapkan pada vektor basis pertama $\begin{pmatrix}  1  \\ -1\\  0   \end{pmatrix}$, hasilnya adalah $\begin{pmatrix}  { \frac{   1  }{  \sqrt{6}  } }  \\ - { \frac{   1  }{  \sqrt{6}  } } \\  0   \end{pmatrix}$, jadi ini merupakan vektor eigen dengan nilai eigen ${ \frac{   1  }{  \sqrt{6}  } }$. Jika diterapkan pada vektor basis kedua $\begin{pmatrix}  2  \\ 0 \\  -1   \end{pmatrix}$, hasilnya adalah

\mavergleichskettedisp
{\vergleichskette
{ \begin{pmatrix}  { \frac{   7 }{  3 \sqrt{6}  } }  \\ { \frac{   1 }{  3 \sqrt{6 } }}  \\  - { \frac{   4  }{  3 \sqrt{6}  } }   \end{pmatrix}
}
{ =} { - { \frac{   1 }{  3 \sqrt{6 } }} \begin{pmatrix}  1  \\ -1\\  0   \end{pmatrix} +   { \frac{   4  }{  3 \sqrt{6}  } } \begin{pmatrix}  2  \\ 0 \\  -1   \end{pmatrix}
}
{ } {
}
{ } {
}
{ } {
}
}
{}{}{.}
Karena itu, nilai eigen yang lain sama dengan ${ \frac{   4  }{  3 \sqrt{6}  } }$, dan matriks diagonal yang merepresentasikannya adalah

\mathdisp {\begin{pmatrix}  { \frac{   1  }{  \sqrt{6}  } }   &   0  \\  0  &  { \frac{   4  }{  3 \sqrt{6}  } }  \end{pmatrix}} { . }

 }



\inputaufgabeklausurloesung
{6}

{

Buktikan teorema isometri untuk transpor paralel pada hiperpermukaan

\mavergleichskette
{\vergleichskette
{ Y
}
{ \subseteq }{ \R^n
}
{  }{
}
{    }{
}
{   }{
}
}
{}{}{.}

}
{

Untuk membuktikan linearitas, misalkan diberikan

\mavergleichskette
{\vergleichskette
{ v,w
}
{ \in }{ T_PY
}
{  }{
}
{    }{
}
{   }{
}
}
{}{}{}
dan

\mavergleichskette
{\vergleichskette
{ r,s
}
{ \in }{ \R
}
{  }{
}
{    }{
}
{   }{
}
}
{}{}{}.
Misalkan
\mathkor {} {F} {dan} {G} {}
adalah medan vektor paralel sepanjang $\gamma$ yang ditentukan secara unik menurut
Teorema 6.9 (Geometri diferensial (Osnabrück 2023))
dengan

\mavergleichskette
{\vergleichskette
{ F(a)
}
{ = }{ v
}
{  }{
}
{    }{
}
{   }{
}
}
{}{}{}
dan

\mavergleichskette
{\vergleichskette
{ G(a)
}
{ = }{ w
}
{  }{
}
{    }{
}
{   }{
}
}
{}{}{.}
Menurut
Lema 6.7 (Geometri diferensial (Osnabrück 2023)),
\mathl{rF+sG}{} adalah medan vektor paralel dengan

\mavergleichskette
{\vergleichskette
{ (rF+sG)(a)
}
{ = }{ v+ w
}
{  }{
}
{    }{
}
{   }{
}
}
{}{}{.}
Karena keunikan dalam
Teorema 6.9 (Geometri diferensial (Osnabrück 2023)),
\mathl{rF+sG}{} dengan demikian adalah medan vektor paralel untuk vektor tangen
\mathl{rv+sw}{.} Oleh karena itu,

\mavergleichskettedisp
{\vergleichskette
{ \Psi_{\gamma} (rv+sw)
}
{ =} { (rF+sG)(b)
}
{ =} { r F(b) +s G(b)
}
{ =} { r \Psi_{\gamma} (v) + s \Psi_{\gamma} (w)
}
{ } {
}
}
{}{}{.}

Untuk membuktikan kesesuaian dengan hasil kali skalar, misalkan kembali diberikan

\mavergleichskette
{\vergleichskette
{ v,w
}
{ \in }{ T_PY
}
{  }{
}
{    }{
}
{   }{
}
}
{}{}{}
dan misalkan $F,G$ adalah medan-medan vektor paralel yang terkait. Kita peroleh

\mavergleichskettedisp
{\vergleichskette
{ \left\langle  F(t) , G(t)  \right\rangle'
}
{ =} { \left\langle  F'(t) , G(t)  \right\rangle  +  \left\langle  F(t) , G'(t)  \right\rangle
}
{ =} { 0
}
{ } {
}
{ } {
}
}
{}{}{,}
karena $F,G$ tangensial dan $F',G'$ ortogonal terhadap ruang tangen. Jadi
\mathl{\left\langle  F(t) , G(t)  \right\rangle}{} konstan sepanjang lintasan. Oleh karena itu,

\mavergleichskettedisp
{\vergleichskette
{ \left\langle \Psi_\gamma(v)  , \Psi_\gamma(w)  \right\rangle
}
{ =} { \left\langle  F(b)  ,  G(b)   \right\rangle
}
{ =} { \left\langle  F(a)  ,  F(a)   \right\rangle
}
{ =} { \left\langle  v  , w  \right\rangle
}
{ } {
}
}
{}{}{.}
Bijektivitasnya juga jelas.

 }



\inputaufgabeklausurloesung
{0}

{

}
{  }



\inputaufgabeklausurloesung
{6}

{

Misalkan $M$ suatu manifold topologis kompak berdimensi $d$,

\mavergleichskette
{\vergleichskette
{ d
}
{ \geq }{ 1
}
{  }{
}
{    }{
}
{   }{
}
}
{}{}{.}
Tunjukkan bahwa terdapat suatu subhimpunan terbuka terbatas

\mavergleichskette
{\vergleichskette
{ U
}
{ \subseteq }{ \R^d
}
{  }{
}
{    }{
}
{   }{
}
}
{}{}{}
dan suatu pemetaan kontinu surjektif
\maabbdisp {\varphi} { U } { M
} {}
.

}
{

Untuk setiap titik
\mathl{P \in M}{}, pilih suatu lingkungan peta terbuka
\mathl{P \in U_P}{} dengan peta
\maabbdisp {\alpha_P} {U_P } {V_P
} {}
dengan
\mathl{V_P \subseteq \R^d}{.} Dengan beralih ke suatu lingkungan bola dari
\mathl{\alpha_P(P) \in V_P}{} beserta prapetanya, kita dapat mengasumsikan bahwa $V_p$ adalah bola-bola terbuka yang jari-jarinya paling besar
\mathl{{ \frac{   1 }{  3 } }}{}. Himpunan-himpunan

\mathbed {U_P} {}
 {P\in M} {}
 {} {} {} {,}
menutupi manifold tersebut. Karena kekompakan $M$, terdapat berhingga banyak titik
\mathl{P_1 , \ldots , P_n}{} sedemikian sehingga

\mathbed {U_i=U_{P_i}} {}
 {i=1 , \ldots , n} {}
 {} {} {} {,}
juga menutupi manifold tersebut. Kita tempatkan bola-bola terbuka
\mathl{V_i=V_{P_i}}{} dalam
\zusatzklammer {\anfuehrung{$\R^d$ yang baru}{}} {} {},
dengan titik-titik pusat

\mathdisp {(1,0 , \ldots , 0), (2,0 , \ldots , 0) , \ldots , (n,0 , \ldots , 0)} { . }

Karena syarat jari-jari, bola-bola ini saling lepas. Kita tinjau himpunan terbuka terbatas
\mathl{U= \bigcup_{i=1}^n V_i}{.} Pemetaan-pemetaan peta $\alpha_i$ menghasilkan pemetaan kontinu

\mathdisp {\varphi_i :V_i \stackrel{ \left(  \alpha_i  \right)^{-1} }{\longrightarrow} U_i \longrightarrow M} { . }

Karena saling lepas, pemetaan-pemetaan ini mendefinisikan, melalui
\mathl{\varphi(z)= \varphi_i(z)}{,} apabila
\mathl{z \in V_i}{}, suatu pemetaan kontinu
\maabbdisp {\varphi} { U } {M
} {.}
Karena sifat penutupnya, pemetaan ini surjektif.

 }



\inputaufgabeklausurloesung
{11 (1+3+1+2+4)}

{

Kita tinjau himpunan

\mavergleichskettedisp
{\vergleichskette
{N
}
{ =} { \left\{ A= \begin{pmatrix}  a   &   b   \\  c  & d \end{pmatrix}   \mid   A \text{ nilpoten}         \right\}
}
{ \subseteq} {\R^4
}
{ } {
}
{ } {
}
}
{}{}{}
dari
$2\times 2$-\definitionsverweis {matriks}{}{}
real nilpoten, serta himpunan

\mavergleichskettedisp
{\vergleichskette
{M
}
{ =} { N \setminus \{ \begin{pmatrix}  0   &   0  \\  0  &  0 \end{pmatrix}  \}
}
{ } {
}
{ } {
}
{ } {
}
}
{}{}{.}

a) Apakah $N$ terhubung?

b) Tunjukkan bahwa $M$ suatu
\definitionsverweis {submanifold tertutup}{}{}
dari suatu subhimpunan terbuka

\mavergleichskette
{\vergleichskette
{ G
}
{ \subseteq }{ \R^4
}
{  }{
}
{    }{
}
{   }{
}
}
{}{}{.}

c) Tentukan
\definitionsverweis {dimensi}{}{}
$M$.

d) Apakah $M$ terhubung?

e) Tutupi
\mathl{M}{} dengan peta-peta topologis eksplisit.

}
{

a) Setiap matriks nilpoten
\mathl{\begin{pmatrix}  a   &   b   \\  c  & d \end{pmatrix}}{} dapat dihubungkan dengan matriks nol di dalam himpunan matriks nilpoten melalui lintasan linear
\maabbeledisp {} {[0,1]} {N
} { t } {t \begin{pmatrix}  a   &   b   \\  c  & d \end{pmatrix}
} {,}.
Karena itu, $N$ terhubung-lintasan dan dengan demikian juga terhubung.

b) Suatu
$2 \times 2$-\definitionsverweis {matriks}{}{}
bersifat nilpoten tepat ketika jejak dan determinannya sama dengan $0$. Jadi, himpunan $N$ dari matriks-matriks nilpoten dapat dipandang sebagai

\mathdisp {\left\{  (a,b,c,d)  \mid   a+d = 0 \text{ dan } ad-bc = 0         \right\}} {  }

Kita tinjau pemetaan
\maabbeledisp {} {\R^4} {\R^2
} {(a,b,c,d)} { (a+d, ad-bc)
} {.}
Matriks Jacobi pemetaan ini adalah

\mathdisp {\begin{pmatrix}  1 &  0 &  0 &  1 \\  d  &  -c &  -b &  a  \end{pmatrix}} { . }

Pemetaan ini tidak
\definitionsverweis {reguler}{}{,}
di titik nol, tetapi reguler di setiap titik lain pada serat $N$. Memang, jika
\mathl{P \in M}{}, maka dari

\mavergleichskettedisp
{\vergleichskette
{a^2
}
{ =} {bc
}
{ } {
}
{ } {
}
{ } {
}
}
{}{}{}
dan

\mavergleichskettedisp
{\vergleichskette
{b
}
{ =} { c
}
{ =} { 0
}
{ } {
}
{ } {
}
}
{}{}{}
langsung diperoleh

\mavergleichskettedisp
{\vergleichskette
{a
}
{ =} {b
}
{ =} { c
}
{ =} { d
}
{ =} { 0
}
}
{}{}{.}
Jadi, matriks Jacobi memiliki rank maksimal di titik-titik dalam $M$. Dengan

\mavergleichskettedisp
{\vergleichskette
{G
}
{ =} {\R^4 \setminus \{0\}
}
{ } {
}
{ } {
}
{ } {
}
}
{}{}{}
kita dapat memandang $M$ sebagai serat pemetaan terbatas yang reguler di seluruh serat tersebut. Oleh karena itu, menurut
Teorema 8.2 (Geometri diferensial (Osnabrück 2023)),
$M$ adalah submanifold tertutup dari $G$.

c) Menurut
Teorema 8.2 (Geometri diferensial (Osnabrück 2023)),
dimensi $M$ sama dengan

\mavergleichskette
{\vergleichskette
{ 4-2
}
{ = }{ 2
}
{  }{
}
{    }{
}
{   }{
}
}
{}{}{.}

d) Kita tulis

\mavergleichskettedisp
{\vergleichskette
{M_+
}
{ =} { \left\{ (a,b,c,d) \in M  \mid   b>c        \right\}
}
{ } {
}
{ } {
}
{ } {
}
}
{}{}{}
dan

\mavergleichskettedisp
{\vergleichskette
{M_-
}
{ =} { \left\{ (a,b,c,d) \in M  \mid   b<c        \right\}
}
{ } {
}
{ } {
}
{ } {
}
}
{}{}{,}
Keduanya,
\zusatzklammer {sebagai irisan $M$ dengan himpunan yang diberikan oleh
\mathl{b>c}{} dan terbuka di $\R^4$} {} {},
merupakan himpunan terbuka dalam $M$. Matriks-matriks
\mathkor {} {\begin{pmatrix}  0  &   1  \\  0 &  0 \end{pmatrix}} {dan} {\begin{pmatrix}  0  &   0  \\  1 &  0 \end{pmatrix}} {}
menunjukkan bahwa keduanya tidak kosong. Selain itu, keduanya menutupi seluruh $M$. Memang, jika

\mavergleichskette
{\vergleichskette
{b
}
{ = }{ c
}
{  }{
}
{    }{
}
{   }{
}
}
{}{}{}
maka dari

\mavergleichskettedisp
{\vergleichskette
{ 0
}
{ =} { ad-bc
}
{ =} { -a^2 -b^2
}
{ } {
}
{ } {
}
}
{}{}{}
langsung diperoleh

\mavergleichskette
{\vergleichskette
{a
}
{ = }{b
}
{ = }{ c
}
{ =   }{ d
}
{ =  }{ 0
}
}
{}{}{,}
dan titik tersebut tidak termasuk dalam $M$. Jadi terdapat penutup oleh dua himpunan terbuka tak kosong yang saling lepas; karena itu $M$ tidak terhubung.

e) Kita bekerja dengan pemetaan
\maabbeledisp {\psi} {\R^2 \setminus \{(0,0)\}} { M_+
} {(a,u)} { (a, u + \sqrt{a^2+u^2}, u - \sqrt{a^2 +u^2},-a)
} {.}
Karena

\mavergleichskettealign
{\vergleichskettealign
{ad-bc
}
{ =} { -a^2 - \left(  u + \sqrt{a^2+u^2}  \right) \left( u - \sqrt{a^2+u^2}  \right)
}
{ =} { -a^2 - \left(  u^2 - \left(  a^2+u^2 \right)   \right)
}
{ =} { 0
}
{ } {
}
}
{}
{}{}
syarat determinan dipenuhi, dan karena

\mavergleichskettedisp
{\vergleichskette
{b-c
}
{ =} { u + \sqrt{a^2+u^2} - \left(  u- \sqrt{a^2 +u^2})  \right)
}
{ =} { 2 \sqrt{a^2+u^2}
}
{ >} { 0
}
{ } {
}
}
{}{}{}
citranya termasuk dalam $M_+$. Pemetaan
\maabbeledisp {\varphi} {M_+} { \R^2 \setminus \{(0,0)\}
} {(a,b,c,d)} { \left(  a  , \,   { \frac{  b+c }{  2 } }                   \right)
} {,}
merupakan invers dari pemetaan yang diberikan. Di sini,

\mavergleichskettedisp
{\vergleichskette
{ \varphi \circ \psi
}
{ =} {
\operatorname{Id}_{ \R^2 \setminus \{(0,0)\}  }
}
{ } {
}
{ } {
}
{ } {
}
}
{}{}{}
jelas. Identitas yang lain diperoleh dari

\mavergleichskettealign
{\vergleichskettealign
{ { \frac{  b+c }{  2 } } + \sqrt{a^2 + \left( { \frac{  b+c }{  2 } }  \right)^2 }
}
{ =} { { \frac{  b+c }{  2 } } + { \frac{   1 }{  2 } }  \sqrt{4 a^2 + b^2 +c^2 +2bc }
}
{ =} { { \frac{  b+c }{  2 } } + { \frac{   1 }{  2 } }  \sqrt{ b^2 +c^2 -2bc }
}
{ =} { { \frac{  b+c }{  2 } } + { \frac{  b-c }{  2 } }
}
{ =} {b
}
}
{}
{}{}
dan

\mavergleichskettedisp
{\vergleichskette
{ { \frac{  b+c }{  2 } } - \sqrt{a^2 + \left(   { \frac{  b+c }{  2 } }   \right)^2 }
}
{ =} { { \frac{  b+c }{  2 } } - { \frac{  b-c }{  2 } }
}
{ =} { c
}
{ } {}
{ } {}
}
{}{}{.}
Kedua pemetaan kontinu, sehingga diperoleh suatu
\definitionsverweis {homeomorfisme}{}{}.
Untuk
\mathl{M_-}{}, tukarkan peran
\mathkor {} {b} {dan} {c} {.}

 }



\inputaufgabeklausurloesung
{5}

{

Katakan sesuatu yang cerdas tentang pita Möbius.

}
{  }



\inputaufgabeklausurloesung
{5}

{

Kita tinjau pemetaan
\maabbeledisp {\varphi} { \R^3 } { \R^2
} { (x,y,z) } { \left(  xyz^2,xy-z^3  \right)
} {.}
Hitung matriks pemetaan
\maabbdisp {\bigwedge^2 T_P (\varphi)} { \bigwedge^2 T_P \R^3 } { \bigwedge^2 T_{\varphi(P)}  \R^2
} {}
di titik

\mavergleichskette
{\vergleichskette
{ P
}
{ = }{ (1,3,5)
}
{  }{
}
{    }{
}
{   }{
}
}
{}{}{}
terhadap suatu basis yang sesuai.

}
{

Secara umum, matriks Jacobi dari $\varphi$ adalah

\mathdisp {\begin{pmatrix}  yz^2 &  xz^2 &  2xyz \\  y  &  x  &  -3z^2 \end{pmatrix}} { . }

Untuk titik

\mavergleichskette
{\vergleichskette
{ P
}
{ = }{ (1,3,5)
}
{  }{
}
{    }{
}
{   }{
}
}
{}{}{}
matriks Jacobinya adalah

\mathdisp {\begin{pmatrix}  75 &  25 &  30 \\  3 &  1 &  -75  \end{pmatrix}} {  }

Matriks ini merepresentasikan pemetaan linear
\maabbdisp {L} {T_P\R^3 \cong \R^3} { T_{\varphi(P)} \R^2 \cong \R^2
} {}
terhadap basis-basis standar. Kita tentukan representasi matriks untuk pemetaan
\maabbdisp {\bigwedge^2 L} {\bigwedge^2 \R^3 } { \bigwedge^2 \R^2
} {}
terhadap basis
$e_1 \wedge e_2,\, e_1 \wedge e_3$ , $e_2 \wedge e_3$
\zusatzklammer {di ruas kiri} {} {}
dan
\mathl{e_1 \wedge e_2}{}
\zusatzklammer {di ruas kanan} {} {.}
Untuk itu, kita harus menghitung citra hasil kali-hasil kali baji tersebut. Kita peroleh

\mavergleichskettealign
{\vergleichskettealign
{ (\bigwedge^2 L ) (e_1 \wedge e_2)
}
{ =} { L(e_1) \wedge L(e_2)
}
{ =} { (75e_1 + 3e_2) \wedge (  25 e_1 + 1 e_2 )
}
{ =} { 75 e_1 \wedge e_2 + 75 e_2 \wedge  e_1
}
{ =} { 75 e_1 \wedge e_2 - 75 e_1 \wedge  e_2
}
}
{
\vergleichskettefortsetzungalign
{ =} { 0
}
{ } {}
{ } {}
{ } {}
}
{}{,}

\mavergleichskettealign
{\vergleichskettealign
{ (\bigwedge^2 L ) (e_1 \wedge e_3)
}
{ =} { L(e_1) \wedge L(e_3)
}
{ =} { (75e_1 + 3e_2) \wedge (  30 e_1 - 75 e_2 )
}
{ =} { - 75^2 e_1 \wedge e_2 + 90 e_2 \wedge  e_1
}
{ =} { (-5625 - 90) e_1 \wedge e_2
}
}
{
\vergleichskettefortsetzungalign
{ =} { -5715 e_1 \wedge e_2
}
{ } {}
{ } {}
{ } {}
}
{}{,}
dan

\mavergleichskettealign
{\vergleichskettealign
{ (\bigwedge^2 L ) (e_2 \wedge e_3)
}
{ =} { L(e_2) \wedge L(e_3)
}
{ =} { (  25 e_1 + 1 e_2 )   \wedge ( 30 e_1 - 75 e_2)
}
{ =} { - 1875 e_1 \wedge e_2 + 30 e_2 \wedge  e_1
}
{ =} { -  1905 e_1 \wedge  e_2
}
}
{}
{}{.}
Jadi, matriks yang merepresentasikannya adalah

\mathdisp {\left(  0,-5715,-1905                     \right)} { . }

 }



\inputaufgabeklausurloesung
{3}

{

Misalkan

\mavergleichskette
{\vergleichskette
{ W
}
{ \subseteq }{ \R^m
}
{  }{
}
{    }{
}
{   }{
}
}
{}{}{}
dan

\mavergleichskette
{\vergleichskette
{ U
}
{ \subseteq }{ \R^n
}
{  }{
}
{    }{
}
{   }{
}
}
{}{}{}
merupakan
\definitionsverweis {subhimpunan-subhimpunan terbuka}{}{}
dan misalkan
\maabbdisp {\psi} { W } { U
} {}
suatu
\definitionsverweis {pemetaan}{}{}
yang
\definitionsverweis {terdiferensialkan secara kontinu}{}{.}
Misalkan
\maabbdisp {f} { U } {\R
} {}
suatu fungsi terdiferensialkan secara kontinu. Simpulkan dari
aturan rantai
bahwa

\mavergleichskettedisp
{\vergleichskette
{ d (\psi^*f)
}
{ =} { \psi^*(df)
}
{ } {
}
{ } {
}
{ } {
}
}
{}{}{,}
dengan $\psi^*$ menyatakan
\definitionsverweis {penarikan balik bentuk diferensial}{}{.}

}
{

Misalkan

\mavergleichskette
{\vergleichskette
{ P
}
{ \in }{  W
}
{  }{
}
{    }{
}
{   }{
}
}
{}{}{}
dan

\mavergleichskette
{\vergleichskette
{ v
}
{ \in }{  \R^m
}
{  }{
}
{    }{
}
{   }{
}
}
{}{}{.}
Di satu pihak,

\mavergleichskettedisp
{\vergleichskette
{ d( \psi^* f) (P,v)
}
{ =} { D_P( f \circ \psi) (v)
}
{ =} { \left(  \left(  D_{\psi(P)} f  \right)  \circ \left(  D_P  \psi  \right)  \right) (v)
}
{ =} { D_{\psi(P)} (f) \left(  D_P \psi(v)  \right)
}
{ } {
}
}
{}{}{.}
Di pihak lain, juga berlaku

\mavergleichskettedisp
{\vergleichskette
{ \left(  \psi^*(df)  \right) (P,v)
}
{ =} { (df) \left(  \psi(P), (D_P \psi )(v)  \right)
}
{ =} { D_{\psi(P)} (f) \left(  D_P \psi (v)  \right)
}
{ } {
}
{ } {
}
}
{}{}{.}

 }



\inputaufgabeklausurloesung
{7}

{

Misalkan $M$ suatu
\definitionsverweis {manifold diferensiabel}{}{}
dengan
\definitionsverweis {basis terhitung bagi topologinya}{}{}
dan misalkan $\omega$ suatu
\definitionsverweis {bentuk volume positif}{}{}
pada $M$. Misalkan
\maabbdisp {\varphi} { L } { M
} {}
suatu
\definitionsverweis {difeomorfisme}{}{}
dengan manifold $L$ dan

\mavergleichskette
{\vergleichskette
{ S
}
{ \subseteq }{ L
}
{  }{
}
{    }{
}
{   }{
}
}
{}{}{}
suatu
\definitionsverweis {subhimpunan terukur}{}{.}
Tunjukkan bahwa

\mavergleichskettedisp
{\vergleichskette
{ \int_S \varphi^* \omega
}
{ =} { \int_{ \varphi(S)} \omega
}
{ } {
}
{ } {
}
{ } {
}
}
{}{}{.}

}
{

Menurut
Lema 15.2 (Geometri diferensial (Osnabrück 2023)),
kita dapat mengasumsikan bahwa
\mathkor {} {S} {dan} {\varphi(S)} {}
masing-masing seluruhnya terletak dalam satu peta dari
\mathkor {} {L} {dan} {M} {},
katakanlah

\mavergleichskette
{\vergleichskette
{ S
}
{ \subseteq }{ U
}
{  }{
}
{    }{
}
{   }{
}
}
{}{}{}
dengan peta
\maabbdisp {\alpha} { U } { U' \subseteq \R^n
} {}
dan

\mavergleichskette
{\vergleichskette
{ \varphi(S)
}
{ \subseteq }{ V
}
{  }{
}
{    }{
}
{   }{
}
}
{}{}{}
dengan peta
\maabbdisp {\beta} { V } { V' \subseteq \R^n
} {.}
Dengan memperkecil peta-peta itu lebih lanjut, kita dapat mengasumsikan bahwa
\mathkor {} {U} {dan} {V} {}
difeomorfik satu sama lain melalui $\varphi$. Dengan demikian terdapat diagram komutatif difeomorfisme

\mathdisp {\begin{matrix}  U  & \stackrel{ \varphi }{\longrightarrow} &  V  &  \\ \!\!\!\!\! \alpha \downarrow & & \downarrow \beta \!\!\!\!\!  & \\  U'  & \stackrel{ \beta \circ \varphi \circ \alpha^{-1} }{\longrightarrow} &  V' & \! \! \! \!\!\!\!\!   \\ \end{matrix}} {  }

yang menginduksi diagram komutatif

\mathdisp {\begin{matrix}  S  & \stackrel{ \varphi }{\longrightarrow} &  \varphi (S)  &  \\ \!\!\!\!\! \alpha \downarrow & & \downarrow \beta \!\!\!\!\!  & \\  \alpha (S) & \stackrel{ \beta \circ \varphi \circ \alpha^{-1} }{\longrightarrow} &  \beta ( \varphi (S) ) & \! \! \! \!\!\!\!\!   \\ \end{matrix}} {  }

atas himpunan-himpunan terukur. Untuk bentuk volume $\omega$ pada $V$ berlaku

\mavergleichskettedisp
{\vergleichskette
{ \alpha^{ {-1}*} ( \varphi^* \omega)
}
{ =} { \left(  \beta \circ \varphi \circ \alpha^{-1}  \right)^*  ( \beta^{ {-1}*} \omega)
}
{ } {
}
{ } {
}
{ } {
}
}
{}{}{.}
menurut
Lema 14.7 (Geometri diferensial (Osnabrück 2023))  (4).
Bentuk volume $\beta^{ {-1}*} \omega$ pada $V'$ memiliki representasi
\mathl{g dy_1  \wedge \ldots \wedge dy_n}{} dengan suatu fungsi terukur
\maabb {g} {V'} {\R
} {},
dan menurut definisi, $\int_{\varphi(S)} \omega$ sama dengan $\int_{\beta( \varphi(S))} g d \lambda^n$. Demikian pula,
\mathl{\alpha^{ {-1}*} ( \varphi^* \omega)}{} pada $U'$ memiliki representasi berbentuk
\mathl{f dx_1 \wedge \ldots \wedge dx_n}{.} Bentuk volume pada $U'$ ini sama dengan tarikan balik dari
\mathl{g dy_1  \wedge \ldots \wedge dy_n}{}, dan menurut
Korolari 14.9 (Geometri diferensial (Osnabrück 2023)),
tarikan balik tersebut sama dengan

\mathdisp {(g \circ \psi) \cdot \det  \left(  \left(  { \frac{   \partial   \psi_i   }{  \partial   x_j   } }  \right)_{1 \leq  i, j \leq n}  \right) dx_1 \wedge \ldots \wedge dx_n} {  }

\zusatzklammer {dengan

\mavergleichskettek
{\vergleichskettek
{ \psi
}
{ = }{ \beta \circ \varphi \circ \alpha^{-1}
}
{  }{
}
{    }{
}
{   }{
}
}
{}{}{}} {} {.}
Dengan demikian, pernyataan tersebut mengikuti dari
Ujian dengan solusi (Teori ukuran dan integrasi (Osnabrück 2022-2023)) (Teori ukuran dan integrasi (Osnabrück 2022-2023))
(Geometri diferensial (Osnabrück 2023)).

 }



\inputaufgabeklausurloesung
{2}

{

Pada
\definitionsverweis {bundel vektor trivial}{}{}
\maabbdisp {p} {\R^2 \times \R} { \R^2
} {}
ber-
\definitionsverweis {rank}{}{}
$1$ di atas $\R^2$, kita tinjau
\definitionsverweis {koneksi linear}{}{,}
yang diberikan oleh
\definitionsverweis {simbol-simbol Christoffel}{}{}

\mavergleichskette
{\vergleichskette
{ \Gamma_1 (x,y)
}
{ = }{ x^5-x^2y^2+3y^4
}
{  }{
}
{    }{
}
{   }{
}
}
{}{}{}
dan

\mavergleichskette
{\vergleichskette
{ \Gamma_2 (x,y)
}
{ = }{ 7x^3 +xy^2-4y^5
}
{  }{
}
{    }{
}
{   }{
}
}
{}{}{.}
Hitung
\definitionsverweis {operator kelengkungan}{}{}
$R \left(  \partial_1, \partial_2  \right)$.

}
{

Menurut
Fakta *****,
berlaku

\mavergleichskettealign
{\vergleichskettealign
{ R \left(  \partial_1, \partial_2  \right)
}
{ =} { \partial_1 \Gamma_2 - \partial_2 \Gamma_1
}
{ =} { \partial_1 \left(  7x^3 +xy^2-4y^5  \right) -  \partial_2 \left(   x^5-x^2y^2+3y^4   \right)
}
{ =} {  21x^2 +y^2 - 5x^4 -2xy^2
}
{ } {
}
}
{}
{}{.}

 }



\inputaufgabeklausurloesung
{0}

{

}
{  }



\inputaufgabeklausurloesung
{6}

{

Buktikan teorema tentang retraksi ke batas pada manifold.

}
{

Menurut
Teorema 21.8 (Geometri diferensial (Osnabrück 2023)),
batas
\mathl{\partial M}{} merupakan
\definitionsverweis {manifold diferensiabel}{}{}
\definitionsverweis {berorientasi}{}{}
\zusatzklammer {tanpa batas} {} {.}
Oleh karena itu, menurut
Teorema 22.11 (Geometri diferensial (Osnabrück 2023)),
terdapat suatu
\definitionsverweis {bentuk volume positif}{}{}
\definitionsverweis {yang terdiferensialkan secara kontinu}{}{}
$\tau$ pada
\mathl{\partial M}{.} Kita mempunyai
\mathl{\int_{  \partial M  }  \tau >0}{.}
\definitionsverweis {Turunan eksterior}{}{}
dari bentuk volume $\tau$ adalah $0$. Andaikan terdapat suatu
\definitionsverweis {pemetaan yang terdiferensialkan secara kontinu}{}{}
\maabbdisp {\varphi} { M } { \partial M
} {}
dengan
\mathl{\varphi \vert_{\partial M} = \operatorname{Id}_{\partial M}}{}. Maka
\definitionsverweis {bentuk tarikan balik}{}{}

\mathl{\varphi^* \tau}{} adalah
$(n-1)$-\definitionsverweis {bentuk diferensial}{}{}
pada $M$ yang pembatasannya pada batas berimpit dengan $\tau$. Dengan menggunakan
Teorema 23.2 (Geometri diferensial (Osnabrück 2023))
dan
Teorema 20.4 (Geometri diferensial (Osnabrück 2023))  (5),
kita memperoleh

\mavergleichskettealign
{\vergleichskettealign
{
\int_{  \partial M  }  \tau
}
{ =} {
\int_{  \partial  M  }  \varphi^* \tau
}
{ =} {
\int_{   M  }  d(\varphi^* \tau)
}
{ =} {
\int_{   M  }  \varphi^* (d\tau)
}
{ =} {
\int_{   M  }  \varphi^* (0)
}
}
{
\vergleichskettefortsetzungalign
{ =} { 0
}
{ } {}
{ } {}
{ } {}
}
{}{.}
Ini merupakan suatu kontradiksi.

 }
